% !TEX root = ../QuantumComputing.tex
%% Lecture 1: Quantum Circuits and the Circuit Model
%% Based on original notes by T.J. Osborne, enriched and unified

\section{Quantum gate operations and the circuit model}\label{sec:circuits}

\begin{historical}[From classical to quantum computation]
The idea that quantum mechanics could be harnessed for computation was first
proposed by Richard Feynman in 1982~\cite{Feynman1982}, who observed that
simulating quantum systems on classical computers appears to require
exponential resources.  David Deutsch formalised the notion of a quantum
Turing machine in 1985, and the quantum circuit model emerged as the standard
framework for designing quantum algorithms. Today, the circuit model
underpins essentially all practical quantum computing architectures, from
superconducting qubits to trapped ions.
\end{historical}

\subsection{Classical computation}\label{subsec:classical}

\subsubsection{Boolean circuits}

We begin with the theoretical model underlying classical computation.

A \emph{boolean circuit} is made up of \emph{wires} and \emph{gates} which
carry information around and perform simple computational tasks,
respectively.  These are theoretical abstractions and may, or may not, be
directly identifiable with actual wires and gates (made from transistors).

\begin{center}
\begin{tikzpicture}
  % NOT gate circuit diagram
  \draw[qwire] (0,0) -- (1.5,0);
  \node[gate] at (2,0) {$\NOT$};
  \draw[qwire] (2.5,0) -- (4,0);
  \node[anchor=east,font=\small] at (0,0) {$a$};
  \node[anchor=west,font=\small] at (4,0) {$\overline{a}$};
\end{tikzpicture}
\end{center}
This example shows a simple circuit which takes as input a single bit, $a$.
This bit is passed through a $\NOT$ gate, which flips the bit, taking $1$ to $0$
and $0$ to $1$. The wires before and after the $\NOT$ gate serve merely to
carry the bit to and from the gate; they represent movement of the bit
through space, or perhaps just through time.

More generally, a \emph{circuit} may involve many input and output bits, many
wires, and many logical gates. A \emph{logic gate} is a \emph{boolean
  function}
\begin{equation}
  \eqbox{f:\{0,1\}^{\times k} \rightarrow \{0,1\}^{\times l}}
\end{equation}
from some fixed number $k$ of input bits to some fixed number $l$ of output
bits.  For example, the $\NOT$ gate is a gate with one input bit and one
output bit which computes the function $f(a) = 1\oplus a \equiv \overline{a}$,
where $a$ is a single bit, and $\oplus$ is modulo~$2$ addition.  It is also
usual to make the convention that no loops are allowed in the circuit, to
avoid possible instabilities. We say such a circuit is \emph{acyclic} and we
adhere to the convention that circuits in the circuit model of computation be
acyclic.

There are many other elementary logic gates which are useful for computation.
A partial list includes the $\AND$ gate, the $\OR$ gate, the $\XOR$ gate,
the $\NAND$ gate, and the $\NOR$ gate.  Each of these gates takes two bits
as input, and produces a single bit as output:

\begin{definition}[Standard logic gates]\label{def:logic-gates}
  The standard two-input, one-output logic gates are defined by their truth
  tables:
  \begin{center}
    \begin{tabular}{cc|ccccc}
      \toprule
      $a$ & $b$ & $\AND$ & $\OR$ & $\XOR$ & $\NAND$ & $\NOR$ \\
      \midrule
      0 & 0 & 0 & 0 & 0 & 1 & 1 \\
      0 & 1 & 0 & 1 & 1 & 1 & 0 \\
      1 & 0 & 0 & 1 & 1 & 1 & 0 \\
      1 & 1 & 1 & 1 & 0 & 0 & 0 \\
      \bottomrule
    \end{tabular}
  \end{center}
\end{definition}

There are two important ``gates'' missing from this table, namely the
$\textsc{fanout}$ gate and the $\textsc{crossover}$ gate: In circuits we
often allow bits to ``divide'', replacing a bit with two copies of itself, an
operation referred to as $\textsc{fanout}$.  We also allow bits to
$\textsc{crossover}$, that is, the value of two bits are interchanged. A
third operation missing, not really a logic gate at all, is to allow the
preparation of extra \emph{ancilla} or \emph{work bits}, to allow extra
working space during the computation.

\begin{remark}[No-cloning in quantum mechanics]
  In classical computing, $\textsc{fanout}$ (copying a bit) is
  straightforward.  In quantum mechanics, however, the \emph{no-cloning
    theorem} states that there is no unitary operation mapping
  $\ket{\psi}\ket{0}\mapsto\ket{\psi}\ket{\psi}$ for all states
  $\ket{\psi}$.  The proof is elementary: if such a unitary $U$ existed,
  then for two states $\ket{\phi}$, $\ket{\psi}$ we would have
  $U\ket{\phi}\ket{0} = \ket{\phi}\ket{\phi}$ and
  $U\ket{\psi}\ket{0} = \ket{\psi}\ket{\psi}$.  Taking the inner product:
  $\braket{\phi}{\psi} = \braket{\phi}{\psi}^2$, which forces
  $\braket{\phi}{\psi}\in\{0,1\}$---so $U$ can only clone orthogonal or
  identical states, not arbitrary ones.
  The impossibility of cloning is a fundamental constraint that distinguishes
  quantum circuits from classical ones.
\end{remark}

\subsubsection{Universality of classical gates}

Just a few fixed gates can be used to compute any function
$f:\{0,1\}^{\times n} \rightarrow \{0,1\}^{\times m}$ whatsoever.  One can
prove this for the simplified case of a function
$f:\{0,1\}^{\times n} \rightarrow \{0,1\}$ with $n$ input bits and a single
output bit.  The general universality proof follows immediately from this
special case.

\begin{theorem}[Universality of $\AND$, $\XOR$, $\NOT$]\label{thm:classical-universality}
  The gate set $\{\AND, \XOR, \NOT\}$, together with wires, ancilla bits,
  $\textsc{fanout}$, and $\textsc{crossover}$, is universal for classical
  computation. That is, any boolean function
  $f:\{0,1\}^{\times n}\to\{0,1\}$ can be computed by a circuit built from
  these elements.
\end{theorem}

\begin{intuition}[Why a small gate set suffices]
  The proof proceeds by induction on $n$.  For $n=1$ there are only four
  boolean functions $\{0,1\}\to\{0,1\}$, each implementable directly. For
  the inductive step, any $(n+1)$-bit function $f$ can be decomposed as
  \[
    f(x_0, x_1,\ldots,x_n) =
    \overline{x_0}\cdot f_0(x_1,\ldots,x_n)
    \;\oplus\;
    x_0 \cdot f_1(x_1,\ldots,x_n),
  \]
  where $f_0$ and $f_1$ are the restrictions to $x_0=0$ and $x_0=1$
  respectively. By induction, both $f_0$ and $f_1$ have circuits, and
  combining them with a multiplexer built from $\AND$, $\XOR$, and $\NOT$
  completes the construction.
\end{intuition}

\begin{remark}
  The $\NAND$ gate alone is universal: it can simulate $\NOT$ (by feeding the same input to both ports), $\AND$ (by applying $\NOT$ to the output of $\NAND$), and $\XOR$ (via a standard construction).
\end{remark}


\subsection{The quantum circuit model}\label{subsec:qcircuit-model}

For us the term \emph{quantum computer} is synonymous with the \emph{quantum
  circuit model} of computation.  The key elements of the quantum circuit
model are:

\begin{keyresult}[Elements of the quantum circuit model]
  \begin{enumerate}
    \item \textbf{Classical resources}: a quantum computer consists of a
      \emph{classical} part and a \emph{quantum} part.  In practice, certain
      tasks may be made much easier if parts of the computation can be done
      classically.
    \item \textbf{A suitable state space}: a quantum circuit operates on $n$
      qubits.  The state space is the $2^n$-dimensional complex Hilbert
      space $(\C^2)^{\tp n}$.  Product states of the form
      $\ket{x_1}\ket{x_2}\cdots \ket{x_n}$, where $x_j \in \{0,1\}$, are
      known as \emph{computational basis states}.
    \item \textbf{Ability to prepare states in the computational basis}: any
      computational basis state
      $\ket{x_1}\ket{x_2}\cdots \ket{x_n}$ can be prepared in at most $n$
      steps.
    \item \textbf{Ability to perform quantum gates}: unitary operations
      called \emph{quantum gates} can be applied to any subset of the qubits
      as desired, and a \emph{universal family} of gates can be implemented.
    \item \textbf{Ability to perform measurements in the computational
        basis}: measurements may be performed in the computational basis of
      one or more of the qubits.
  \end{enumerate}
\end{keyresult}

\subsubsection{Single-qubit operations}\label{subsubsec:single-qubit}

Any $2\times 2$ unitary operator is called a \emph{single-qubit operation}.
While there is an uncountable infinity of such operations, quantum circuits
often comprise one of a handful of specific gates.  The most important
single-qubit gates are the \emph{Pauli gates} and the \emph{Hadamard gate}:

\begin{definition}[Pauli gates and Hadamard]\label{def:pauli-hadamard}
  The Pauli gates are
  \[
    X = \begin{pmatrix} 0 & 1 \\ 1 & 0 \end{pmatrix},\quad
    Y = \begin{pmatrix} 0 & -i \\ i & 0 \end{pmatrix},\quad
    Z = \begin{pmatrix} 1 & 0 \\ 0 & -1 \end{pmatrix}.
  \]
  The Hadamard gate is
  \[
    H = \frac{1}{\sqrt{2}}\begin{pmatrix} 1 & 1 \\ 1 & -1 \end{pmatrix}.
  \]
  The phase gates $S$ and $T$ are
  \[
    S = \begin{pmatrix} 1 & 0 \\ 0 & i \end{pmatrix},\quad
    T = \begin{pmatrix} 1 & 0 \\ 0 & e^{i\pi/4} \end{pmatrix}.
  \]
\end{definition}

These gates can be visualised via \emph{quantum circuit diagrams}:
\begin{center}
\begin{tikzpicture}
  % Wire
  \draw[qwire] (0,0) -- (2,0);
  \node[gate] at (1,0) {$U$};
  \node[anchor=east,font=\small] at (0,0) {$\ket{\psi}$};
  \node[anchor=west,font=\small] at (2,0) {$U\ket{\psi}$};
  %
  \begin{scope}[xshift=4cm]
    \draw[qwire] (0,0) -- (5,0);
    \node[gate] at (1,0) {$H$};
    \node[gate] at (2.5,0) {$T$};
    \node[gate] at (4,0) {$S$};
    \node[anchor=east,font=\small] at (0,0) {$\ket{\psi}$};
  \end{scope}
\end{tikzpicture}
\end{center}

\begin{intuition}[The Bloch sphere]
  Any single-qubit pure state $\ket{\psi}=\alpha\ket{0}+\beta\ket{1}$ (with
  $|\alpha|^2+|\beta|^2=1$) can be parameterised as
  $\ket{\psi}=\cos(\theta/2)\ket{0}+e^{i\varphi}\sin(\theta/2)\ket{1}$,
  mapping states to points on the unit sphere in $\R^3$: the \emph{Bloch
    sphere}.  The north pole is $\ket{0}$, the south pole is $\ket{1}$, and
  the equator contains states like $\ket{+}$ and $\ket{-}$.  Single-qubit
  unitaries correspond to rotations of this sphere.

  \begin{center}
    \begin{tikzpicture}
      \blochsphere
      % Example state vector
      \draw[-{Stealth[length=6pt,width=4pt]},cgblue,thick]
        (0,0) -- (0.8,1.1) node[above right,font=\small] {$\ket{\psi}$};
      \draw[dashed,spacecadet!50] (0,0) -- (0.8,0);
      \draw[dashed,spacecadet!50] (0.8,0) -- (0.8,1.1);
    \end{tikzpicture}
  \end{center}
\end{intuition}

It turns out that any $\SU{2}$ operator can be approximated by a product of
a finite number of gates from a suitable generating set; this is the content
of the \emph{Solovay-Kitaev theorem}~\cite{SolovayKitaev2006}.

\begin{theorem}[Solovay-Kitaev~\cite{SolovayKitaev2006}]\label{thm:SK}
  If a set of single-qubit quantum gates generates a dense subset of $\SU{2}$,
  then any element of $\SU{2}$ can be approximated to accuracy $\epsilon$
  using a sequence of $\bigO{\log^c(1/\epsilon)}$ gates from the generating
  set, where $c\approx 3.97$.
\end{theorem}

\subsubsection{Controlled operations}\label{subsubsec:controlled}

``If $A$ is true, then do $B$'' is one of the most useful operations in
computing, both classical and quantum.

The prototypical controlled operation is the controlled-$\NOT$ ($\CNOT$).
This gate has two input qubits: the \emph{control qubit} and the
\emph{target} qubit.  Its quantum circuit diagram is
\begin{center}
\begin{tikzpicture}
  % CNOT circuit
  \qcwires{2}{0}{4}
  \qcctrl{2}{1}
  \qccnottarg{2}{2}
  \qcctrlline{2}{1}{2}
  \qclabel{1}{\ket{c}}
  \qclabel{2}{\ket{t}}
  \node[anchor=west,font=\small] at (4,-1) {$\ket{c}$};
  \node[anchor=west,font=\small] at (4,-2) {$\ket{t\oplus c}$};
\end{tikzpicture}
\end{center}
In the computational basis the action of the $\CNOT$ is given by
$\ket{c}\ket{t}\mapsto \ket{c}\ket{t\oplus c}$; that is, if the control
qubit is set to $\ket{1}$ then the target qubit is flipped, otherwise the
target qubit is left alone.  The matrix representation of $\CNOT$ is:
\begin{equation}
  \CNOT \equiv \begin{pmatrix}
    1 & 0 & 0 & 0  \\
    0 & 1 & 0 & 0  \\
    0 & 0 & 0 & 1  \\
    0 & 0 & 1 & 0  \\
  \end{pmatrix}.
\end{equation}

\begin{intuition}[The $\CNOT$ creates entanglement]
  When applied to the state $\ket{+}\ket{0} = \frac{1}{\sqrt{2}}(\ket{0}+\ket{1})\ket{0}$,
  the $\CNOT$ produces the \emph{Bell state}
  $\frac{1}{\sqrt{2}}(\ket{00}+\ket{11})$, one of the maximally entangled
  two-qubit states.  This illustrates the key difference from classical
  computation: quantum gates can generate entanglement, which is the
  essential resource for quantum speedups.
\end{intuition}

Now that we know how to condition on a single qubit, what about conditioning
on multiple qubits?  An example of multiple-qubit conditioning is the
\emph{Toffoli gate}, which flips the third qubit conditioned on the first
two qubits both being set to one:
\begin{center}
\begin{tikzpicture}
  % Toffoli gate
  \qcwires{3}{0}{4}
  \qcctrl{2}{1}
  \qcctrl{2}{2}
  \qccnottarg{2}{3}
  \qcctrlline{2}{1}{3}
  \qclabel{1}{\ket{a}}
  \qclabel{2}{\ket{b}}
  \qclabel{3}{\ket{c}}
  \node[anchor=west,font=\small] at (4,-1) {$\ket{a}$};
  \node[anchor=west,font=\small] at (4,-2) {$\ket{b}$};
  \node[anchor=west,font=\small] at (4,-3) {$\ket{c\oplus ab}$};
\end{tikzpicture}
\end{center}

More generally, suppose we have $n + k$ qubits and $U$ is a $k$-qubit
unitary operator.  Then we define the \emph{controlled operation}
$\CnU{n}{U}$ by
\begin{equation}
  \CnU{n}{U}\ket{x_1}\ket{x_2}\cdots \ket{x_n}\ket{\psi}
  = \ket{x_1}\ket{x_2}\cdots \ket{x_n}\, U^{x_1\cdot x_2\cdot \cdots \cdot x_n}\ket{\psi},
\end{equation}
where the product $x_1\cdot x_2\cdot \cdots \cdot x_n$ in the exponent of $U$
means the Boolean product. That is, $U$ is applied to the last $k$ qubits if
and only if the first $n$ qubits are all equal to $\ket{1}$.

\subsubsection{Universal quantum gates}\label{subsubsec:universal}

A small set of gates (e.g.\ $\AND$, $\XOR$ and $\NOT$) can compute an
arbitrary classical function.  A similar universality result is true for
quantum computation: a set of gates is said to be \emph{universal for
  quantum computation} if any unitary operation may be approximated to
arbitrary accuracy by a quantum circuit involving only those gates.

\begin{keyresult}[Universality of single-qubit gates and $\CNOT$]
  Any unitary operation on $n$ qubits can be approximated to arbitrary
  accuracy using single-qubit operations and $\CNOT$ gates.  This is one of
  the most important results in quantum computation, and justifies the
  experimental focus on implementing such gates.
\end{keyresult}

In particular, the set $\{H, T, \CNOT\}$ is a widely used \emph{discrete}
universal gate set: $H$ and $T$ generate a dense subset of $\SU{2}$ (the
proof uses the irrationality of $\pi^{-1}\arccos(3/\sqrt{8})$), and
combining with $\CNOT$ gives universality for all $n$-qubit unitaries via
the Solovay-Kitaev theorem (Theorem~\ref{thm:SK}).

The proof proceeds in two steps:

\paragraph{Step 1: Two-level unitary gates are universal.}
Consider a unitary matrix $U$ acting on a $d$-dimensional Hilbert space. We
show how $U$ may be decomposed into a product of \emph{two-level} unitary
matrices, i.e., unitary matrices which act non-trivially on at most two
vector components.

The essential idea may be understood by considering $U \in \SU{3}$:
\begin{equation}
  U = \begin{pmatrix}
    a & d & g \\
    b & e & h \\
    c & f & j
  \end{pmatrix}.
\end{equation}
We construct two-level unitaries that successively zero out entries in each
column, analogous to Gaussian elimination but with unitary row operations.

\textbf{Column 1.} If $b\neq 0$, define a two-level unitary $U_1$ acting
on the $(1,2)$ subspace by the $2\times 2$ block
\[
  \widetilde{U}_1 = \frac{1}{\sqrt{|a|^2+|b|^2}}
  \begin{pmatrix} a^* & b^* \\ -b & a \end{pmatrix},
\]
so $U_1 U$ has a zero in position $(2,1)$. (If $b=0$ already, take
$U_1=\I$.) Next, if the $(3,1)$ entry of $U_1 U$ is nonzero, define
$U_2$ acting on the $(1,3)$ subspace to zero it out.  After $U_2 U_1 U$,
the first column has the form $(e^{i\phi},0,0)^T$.  By unitarity, this
forces $U_2 U_1 U$ to be block-diagonal: a phase times the identity on
component~$1$ and a $2\times 2$ unitary on the $(2,3)$ subspace.  A third
two-level unitary $U_3$ acting on the $(2,3)$ subspace inverts this block,
giving $U_3 U_2 U_1 U = \I$ and hence
$U = U_1^\dagger U_2^\dagger U_3^\dagger$.

\textbf{General case.}
For $U\in\SU{d}$, the same strategy reduces column $1$ to
$(1,0,\ldots,0)^T$ using at most $d-1$ two-level unitaries (one for each
sub-diagonal entry).  The result is a block-diagonal matrix with a
$(d-1)\times(d-1)$ unitary block, which is reduced inductively.  The total
count is $(d-1)+(d-2)+\cdots+1 = d(d-1)/2$ two-level unitaries.

\begin{example}[Decomposing a $3\times 3$ permutation matrix]\label{ex:two-level-3x3}
  Consider the cyclic permutation
  \[
    U = \begin{pmatrix} 0 & 0 & 1 \\ 1 & 0 & 0 \\ 0 & 1 & 0 \end{pmatrix},
  \]
  which maps $\ket{1}\to\ket{2}\to\ket{3}\to\ket{1}$.  The first column is
  $(0,1,0)^T$, so $a=0$, $b=1$.  The formula gives
  $\widetilde{U}_1 = \bigl(\begin{smallmatrix} 0 & 1 \\ -1 & 0
  \end{smallmatrix}\bigr)$, hence
  \[
    U_1 = \begin{pmatrix} 0 & 1 & 0 \\ -1 & 0 & 0 \\ 0 & 0 & 1
    \end{pmatrix}, \qquad
    U_1 U = \begin{pmatrix} 1 & 0 & 0 \\ 0 & 0 & -1 \\ 0 & 1 & 0
    \end{pmatrix}.
  \]
  The first column is now $(1,0,0)^T$, so $U_2=\I$.  The remaining
  $2\times 2$ block $\bigl(\begin{smallmatrix} 0 & -1 \\ 1 & 0
  \end{smallmatrix}\bigr)$ is inverted by
  $U_3 = \diag\bigl(1,\,\bigl(\begin{smallmatrix} 0 & 1 \\ -1 & 0
  \end{smallmatrix}\bigr)\bigr)$, giving $U_3 U_1 U = \I$ and hence
  \[
    U = U_1^\dagger\, U_3^\dagger
    = \underbrace{\begin{pmatrix} 0 & -1 & 0 \\ 1 & 0 & 0 \\ 0 & 0 & 1
    \end{pmatrix}}_{\text{acts on $(1,2)$}}
    \underbrace{\begin{pmatrix} 1 & 0 & 0 \\ 0 & 0 & -1 \\ 0 & 1 & 0
    \end{pmatrix}}_{\text{acts on $(2,3)$}},
  \]
  a product of two two-level unitaries.
\end{example}

\paragraph{Step 2: Two-level unitaries from single-qubit and $\CNOT$ gates.}
Any two-level unitary on $n$ qubits can be implemented using single-qubit
and $\CNOT$ gates, via \emph{Gray codes}.

\begin{definition}[Gray code]\label{def:gray-code}
  Suppose we have distinct binary strings
  $\mathbf{s} = s_1 s_2\cdots s_n$ and $\mathbf{t} = t_1 t_2\cdots t_n$.
  A \emph{Gray code} connecting $\mathbf{s}$ and $\mathbf{t}$ is a sequence
  of binary strings, starting with $\mathbf{s}$ and concluding with
  $\mathbf{t}$, such that adjacent members of the list differ in exactly one
  bit.
\end{definition}

For instance, with $n=3$ qubits, $\mathbf{s} = 111$ and
$\mathbf{t} = 001$, we have the Gray code $111 \to 011 \to 001$.

Let $g_1, \ldots, g_m$ be the elements of a Gray code connecting
$\mathbf{s}$ and $\mathbf{t}$, with $g_1 = \mathbf{s}$ and
$g_m = \mathbf{t}$.  Note that $m \le n+1$ since $\mathbf{s}$ and
$\mathbf{t}$ can differ in at most $n$ positions.

Now consider a two-level unitary that acts as the $2\times 2$ block
$\widetilde{U}$ on the subspace $\vspan\{\ket{\mathbf{s}},\ket{\mathbf{t}}\}$
and as the identity elsewhere.  The circuit construction proceeds in three
stages:
\begin{enumerate}
  \item \textbf{Transfer via $\CNOT$s}: For each step $g_i\to g_{i+1}$ in
    the Gray code (for $i=1,\ldots,m-2$), apply a $\CNOT$ controlled on
    the bits of $g_i$ that target the single bit that changes.  This maps
    $\ket{\mathbf{s}}\mapsto \ket{g_{m-1}}$, which differs from
    $\mathbf{t}$ in only one bit (say bit $p$), while leaving all other
    basis states unchanged.
  \item \textbf{Controlled-$\widetilde{U}$}: Apply $\widetilde{U}$ to qubit
    $p$, controlled on the remaining $n-1$ qubits matching the
    corresponding bits of $g_{m-1}$.  This implements $\widetilde{U}$ on
    the subspace $\vspan\{\ket{g_{m-1}},\ket{\mathbf{t}}\}$.
  \item \textbf{Undo transfers}: Reverse the $\CNOT$ sequence from step 1.
\end{enumerate}
Each step uses $\bigO{n}$ elementary gates (the controlled-$\widetilde{U}$
can be decomposed into $\bigO{n}$ single-qubit and $\CNOT$ gates by standard
techniques), so each two-level unitary costs $\bigO{n}$ gates.

\begin{example}[The $\SWAP$ gate from Gray codes]\label{ex:swap-gray}
  The $\SWAP$ gate on $n=2$ qubits exchanges $\ket{01}$ and $\ket{10}$,
  leaving $\ket{00}$ and $\ket{11}$ unchanged---a two-level unitary with
  $\widetilde{U} = X$ on $\vspan\{\ket{01},\ket{10}\}$.  We apply the
  construction above with $\mathbf{s}=01$ and $\mathbf{t}=10$.

  \textbf{Gray code}: $01 \to 11 \to 10$ (flip bit~1, then bit~2), so
  $g_1=01$, $g_2=11$, $g_3=10$.

  \textbf{Transfer} ($g_1\to g_2$): Bit~1 changes.  Apply $\CNOT_{21}$
  (qubit~2 controls qubit~1).  This maps $\ket{01}\to\ket{11}$ while
  $\ket{10}\to\ket{10}$ (unaffected since qubit~2 is $0$).

  \textbf{Controlled-$X$}: Now $g_2=11$ and $g_3=10$ differ only in bit~2.
  Apply $X$ on qubit~2 controlled by qubit~1, i.e., $\CNOT_{12}$.

  \textbf{Undo transfer}: Apply $\CNOT_{21}$ again.

  The full circuit is $\CNOT_{21}\cdot\CNOT_{12}\cdot\CNOT_{21}$:
  \begin{center}
  \begin{tikzpicture}
    \qcwires{2}{0}{8}
    \qclabel{1}{q_1}
    \qclabel{2}{q_2}
    % CNOT_21: q2 controls q1
    \qcctrl{1.5}{2}
    \qccnottarg{1.5}{1}
    \qcctrlline{1.5}{1}{2}
    % CNOT_12: q1 controls q2
    \qcctrl{4}{1}
    \qccnottarg{4}{2}
    \qcctrlline{4}{1}{2}
    % CNOT_21: q2 controls q1
    \qcctrl{6.5}{2}
    \qccnottarg{6.5}{1}
    \qcctrlline{6.5}{1}{2}
  \end{tikzpicture}
  \end{center}
  This recovers the well-known three-$\CNOT$ decomposition of $\SWAP$.
  The reader can verify the action on all four basis states: e.g.,
  $\ket{01}\xrightarrow{\CNOT_{21}}\ket{11}
  \xrightarrow{\CNOT_{12}}\ket{10}
  \xrightarrow{\CNOT_{21}}\ket{10}$,
  and $\ket{10}\xrightarrow{\CNOT_{21}}\ket{10}
  \xrightarrow{\CNOT_{12}}\ket{11}
  \xrightarrow{\CNOT_{21}}\ket{01}$.
\end{example}

By combining both steps: we need $d(d-1)/2$ two-level unitaries (with
$d=2^n$), each costing $\bigO{n}$ gates.  Thus an arbitrary unitary on $n$
qubits can be implemented using $\bigO{n\cdot 4^n}$ single-qubit and $\CNOT$
gates.

\begin{remark}
  The exponential $\bigO{n\cdot 4^n}$ gate count is unavoidable in the
  worst case (a counting argument shows that most $n$-qubit unitaries
  require $\Omega(4^n)$ gates), but for specific unitaries arising in
  algorithms, far more efficient decompositions exist.
\end{remark}
